%% kapa: exam=1 round=21 year=2010 period=1 form=A subject=economics subject_ko=경제원론 seq=37 number=77
%% source: Q-Net 감정평가사 기출문제 (cst00309) artlSeq=5024181 file=extracted/1008/5024181/2010년 제21회 감정평가사 1차시험 기출문제/1교시.hwp
%% conversion: kordoc 4.24.51 -> tools/gen1.py (auto); review=proofread-sonnet; flags=none
\begin{question}{77}
구축효과(crowding-out effect)에 관한 설명으로 옳은 것은?
\begin{choices}
  \item 조세를 인하할 경우 총수요 증가효과가 이자율 상승으로 인한 민간투자의 감소로 상쇄되는 현상을 의미한다.
  \item 조세를 인상할 경우 민간부문의 가처분소득이 감소하여 소비지출이 감소하는 현상을 의미한다.
  \item 정부지출을 확대할 경우 민간부문의 투자지출이 증가하는 현상을 의미한다.
  \item 통화공급량을 감소시킬 경우 이자율이 상승하여 민간부문의 투자지출이 감소하는 현상을 의미한다.
  \item 자국 환율을 평가절하시킬 경우 해외제품 수입증가로 인하여 국내제품에 대한 수요가 감소하는 현상을 의미한다.
\end{choices}
\end{question}
\answer{1}
